\chapter{Spectral normalization калибровочных связей}

Family selector не требуется для относительной gauge normalization: три
поколения и particle--antiparticle doubling умножают все traces на общий
factor, который сокращается в отношениях.

\section{Следы одного generation atom}

Для физического гиперзаряда предыдущих глав fermionic traces равны
\begin{align}
C_Y&=\sum_{\rm Weyl}d_3d_2Y^2=\frac{10}{3},\\
C_2&=\sum_{\rm Weyl}d_3T_2(R)=2,\\
C_3&=\sum_{\rm Weyl}d_2T_3(R)=2.
\end{align}
После GUT normalization
\begin{equation}
T_1=\sqrt{\frac35}\,Y
\end{equation}
получаем
\begin{equation}
C_1=\frac35C_Y=2=C_2=C_3.
\end{equation}

\begin{theorem}[Spectral gauge relation]
Один общий bosonic spectral coefficient перед тремя gauge kinetic terms
фиксирует на matching scale
\begin{equation}
\boxed{g_1=g_2=g_3,}
\end{equation}
где \(g_1=\sqrt{5/3}\,g_Y\). В physical hypercharge convention это
эквивалентно
\begin{equation}
g_Y^2=\frac35g_2^2,
\qquad
g_2^2=g_3^2,
\qquad
\sin^2\theta_W=\frac38
\end{equation}
на matching scale.
\end{theorem}

Это строгий relative-normalization result стандартного почти-коммутативного
baseline. Он не определяет matching scale.

\section{Замороженные экспериментальные входы}

Для one-loop stress test фиксируются \(\overline{\mathrm{MS}}\)-входы на
\(M_Z=91.1876\,\mathrm{GeV}\):
\begin{equation}
\widehat\alpha^{-1}(M_Z)=127.955,
\qquad
\sin^2\widehat\theta_W(M_Z)=0.23122,
\qquad
\alpha_s(M_Z)=0.1180.
\end{equation}
Последние два значения и uncertainties \(0.00006\), \(0.0009\) взяты из
PDG 2026; для running electromagnetic coupling используется стандартный
benchmark uncertainty \(0.010\).

Отсюда
\begin{equation}
g_1(M_Z)=0.46142,
\qquad
g_2(M_Z)=0.65172,
\qquad
g_3(M_Z)=1.21772.
\end{equation}

\section{Minimal one-loop RG gate}

Для Standard Model с тремя поколениями и одним Higgs doublet
\begin{equation}
(b_1,b_2,b_3)=\left(\frac{41}{10},-\frac{19}{6},-7\right),
\end{equation}
а
\begin{equation}
\frac1{g_i^2(\mu)}=
\frac1{g_i^2(M_Z)}-\frac{b_i}{8\pi^2}\log\frac{\mu}{M_Z}.
\end{equation}
Pairwise equality возникает при различных scales:
\begin{align}
g_1=g_2:&\quad \mu_{12}=1.03\times10^{13}\,\mathrm{GeV},\\
g_1=g_3:&\quad \mu_{13}=2.43\times10^{14}\,\mathrm{GeV},\\
g_2=g_3:&\quad \mu_{23}=9.73\times10^{16}\,\mathrm{GeV}.
\end{align}
Следовательно, общего triple intersection нет.

На наиболее близком pairwise crossing \(g_1=g_3\) получаем
\begin{equation}
g_1=g_3=0.55806,
\qquad
g_2=0.53438,
\end{equation}
то есть относительный deficit \(g_2\) составляет примерно \(4.24\%\).
Corner scan экспериментальных uncertainties оставляет его в узком диапазоне
и не снимает mismatch.

\begin{theorem}[Minimal spectral-unification no-go]
Точная spectral trace relation доказана на matching scale, но минимальный
Standard Model без threshold corrections и новых states не переносит её к
измеренным couplings: one-loop trajectories не имеют общего пересечения.
\end{theorem}

\begin{nogocriterion}[Запрет свободного threshold repair]
Нельзя вводить independent threshold shifts для трёх gauge factors после
просмотра mismatch. Следующий repair допустим только если masses и
representations новых states уже следуют из конечной геометрии и одновременно
проходят другой sector gate.
\end{nogocriterion}

Машинный trace и RG ledger сохранён в
\path{s2t_v4_spectral_gauge_normalization_gate_results.json}.